\section{Introduction}
\label{sec:intro}

Multi-site statistical analysis is increasingly constrained by data that
cannot be centralised: electronic health records behind institutional
firewalls, trials governed by separate data-use agreements, claims
databases that no single sponsor controls. Federated learning
\citep{mcmahan2017,kairouz2021} addresses this by exchanging model
parameters rather than records. A central server broadcasts a current
parameter vector, each site improves it on its own data for a while, and
the server aggregates what comes back.

The phrase ``for a while'' hides the design parameter this article is
about. Let $E$ denote the number of local passes---here, full
coordinate-descent (CD) epochs---each site performs between aggregation
rounds. Increasing $E$ reduces the number of communication rounds, which is
the resource federated protocols are usually built to conserve. It also
lets each site travel further toward its own optimum before being pulled
back, which is client drift \citep{li2020convergence,karimireddy2020}. For
smooth losses the consequence of drift is a worse fit. For the Lasso
\citep{tibshirani1996} the consequence is sharper, because the
$\ell_1$ penalty sets coordinates exactly to zero: drift can lead sites to
prune different predictors, and the damage then appears as a different
variable list, not merely a larger residual. When the deliverable of the
analysis is a candidate biomarker panel or a set of covariates for a safety
model, that difference is the whole result.

This article studies the choice of $E$ for a synchronous federated
coordinate-descent Lasso in which each site tunes its own penalty on a
local validation split and the server aggregates by sample-size weight.
Our starting point was an empirical observation that is easy to reproduce
and easy to misread: as $E$ grows, the active set of the aggregated
estimator shrinks and its agreement with a centralised fit deteriorates,
which invites the conclusion that local epochs damage variable selection
through drift. The main contribution of this article is to show that this
reading is, in large part, wrong, and to identify what the aggregation rule
is actually doing.

\paragraph{Contributions.} Our contributions are as follows.

\begin{enumerate}
\item \textbf{The aggregated iterate is not sparse.} Averaging is a linear
operation and soft-thresholding is not; a convex combination of sparse
vectors carries the \emph{union} of their supports. We prove
(Proposition~\ref{prop:orth}) that under orthogonal site designs the
federated iteration reaches its fixed point after a single round
\emph{regardless of $E$}, that this fixed point is
$\sum_j w_j \soft(z_j, \lambda_j)$, and that its active set strictly
contains the active set of the minimiser of the aggregated objective
whenever the sites disagree. Any dependence on $E$ in the orthogonal case
is therefore exactly zero, and all observed $E$-dependence must come from
correlation among predictors.

\item \textbf{The implied objective is the pooled Lasso, and the algorithm
does not minimise it.} With sample-size weights the aggregated objective
collapses exactly to the pooled Lasso objective at the weighted-average
penalty $\bar\lambda$ (Lemma~\ref{lem:pooled}). We characterise the fixed
point of the general iteration as an averaged-stationarity condition
(Proposition~\ref{prop:fixed}) that does not coincide with stationarity of
that objective, and we quantify the resulting gap empirically against
consensus ADMM \citep{boyd2011}, which minimises the same objective
exactly.

\item \textbf{A cheap remedy, and an honest accounting of what it fixes.}
We propose \textsc{FedAvg-ST}, a single additional round in which the
server applies a soft-threshold whose level is selected by federated
validation using only scalar residual summaries. It restores sparsity and a
large share of the selection accuracy, but it does not remove the
optimisation bias; consensus ADMM does, at a higher communication price.

\item \textbf{Communication is not the only budget.} Ranking schedules by
communication rounds alone reverses once local computation is counted:
across our designs, the large-$E$ schedules that transmit least perform the
most total coordinate-descent work.
\end{enumerate}

\paragraph{Related work.} The local-epoch trade-off is well studied for
smooth objectives, where local SGD and its analyses
\citep{stich2019,li2020convergence} and drift corrections such as FedProx
\citep{li2020fedprox} and SCAFFOLD \citep{karimireddy2020} are standard;
composite and proximal federated methods have been developed more recently
\citep{tran2021}. In the statistics literature, communication-efficient
distributed sparse estimation has been approached through one-shot and
surrogate-likelihood constructions \citep{lee2017,battey2018,jordan2019}
and through distributed optimisation with exact consensus
\citep{boyd2011,duchi2012,wang2017}. What has received less attention is
the interaction between parameter averaging and the selection event itself:
the support-recovery literature \citep{zhaoyu2006,meinshausen2010} is
written for a single estimator fitted on pooled data, and the federated
literature typically evaluates loss rather than active sets. The
observation that averaging sparse vectors is not sparse is elementary; its
consequence---that $E$-versus-support comparisons for federated Lasso are
confounded by the aggregation rule---appears not to have been stated
explicitly, and is the gap this article fills.

The remainder of the article is organised as follows.
Section~\ref{sec:setup} fixes notation and states the algorithm.
Section~\ref{sec:theory} contains the analytical results.
Section~\ref{sec:method} introduces \textsc{FedAvg-ST} and the comparators.
Sections~\ref{sec:sim} and~\ref{sec:results} describe and report the Monte
Carlo study, Section~\ref{sec:realdata} an illustration on a three-site
data set, and Section~\ref{sec:discussion} discusses limitations.
